\documentclass[11pt]{article}
\usepackage{mystyle}

\title{Rosen --- \S 2.1: Sets\\ \large Selected Exercises}
\author{Alston}
\date{Semester 2, 2026}

\begin{document}
\maketitle

% ======================================================================
\begin{exercise}{11}
    Determine whether each of these statements is true or false.
    \begin{enumerate}[label=(\alph*)]
        \item $x \in \{x\}$
        \item $\{x\} \subseteq \{x\}$
        \item $\{x\} \in \{x\}$
        \item $\{x\} \in \{\{x\}\}$
        \item $\emptyset \subseteq \{x\}$
        \item $\emptyset \in \{x\}$
    \end{enumerate}
\end{exercise}

% ======================================================================
\begin{exercise}{17}
    Suppose that $A$, $B$, and $C$ are sets such that $A \subseteq B$ and
    $B \subseteq C$. Show that $A \subseteq C$.
\end{exercise}

% ======================================================================
\begin{exercise}{23}
    How many elements does each of these sets have where $a$ and $b$ are
    distinct elements?
    \begin{enumerate}[label=(\alph*)]
        \item $\mathcal{P}(\{a, b, \{a, b\}\})$
        \item $\mathcal{P}(\{\emptyset, a, \{a\}, \{\{a\}\}\})$
        \item $\mathcal{P}(\mathcal{P}(\emptyset))$
    \end{enumerate}
\end{exercise}

% ======================================================================
\begin{exercise}{31}
    Let $A$ be a set. Show that $\emptyset \times A = A \times \emptyset
    = \emptyset$.
\end{exercise}

% ======================================================================
\begin{exercise}{37}
    How many different elements does $A^n$ have when $A$ has $m$
    elements and $n$ is a positive integer?
\end{exercise}

% ======================================================================
\begin{exercise}{45}
    The defining property of an ordered pair is that two ordered pairs
    are equal if and only if their first elements are equal and their
    second elements are equal. Surprisingly, instead of taking the
    ordered pair as a primitive concept, we can construct ordered pairs
    using basic notions from set theory. Show that if we define the
    ordered pair $(a,b)$ to be $\{\{a\},\{a,b\}\}$, then $(a,b) = (c,d)$
    if and only if $a = c$ and $b = d$.
    [\emph{Hint}: First show that $\{\{a\},\{a,b\}\} = \{\{c\},\{c,d\}\}$
    if and only if $a = c$ and $b = d$.]
\end{exercise}

% ======================================================================
\begin{exercise}{46}
    This exercise presents Russell's paradox. Let $S$ be the set that
    contains a set $x$ if the set $x$ does not belong to itself, so that
    $S = \{x \mid x \notin x\}$.
    \begin{enumerate}[label=(\alph*)]
        \item Show the assumption that $S$ is a member of $S$ leads to a
        contradiction.
        \item Show the assumption that $S$ is not a member of $S$ leads
        to a contradiction.
    \end{enumerate}
    By parts (a) and (b) it follows that the set $S$ cannot be defined as
    it was. This paradox can be avoided by restricting the types of
    elements that sets can have.
\end{exercise}

\end{document}